\section{A boundary-value calculus uniform in the itinerary length}
\label{sec:action}
The circular formulas follow from the general boundary-value lemma.
We record their explicit expressions, including the even analytic
improvements, to keep this specialization self-contained at the level
of its calculations. The uniform contraction and determinant argument
are those already proved in Lemma~\ref{lem:g-relative}.

\begin{lemma}[Uniform stationary segment]\label{lem:bridge}
For endpoints $|u|,|v|<a_1$, with $a_1>0$ independent of $j$, the
interior equations of \eqref{eq:action} have a unique small solution.
The solution and $W_{j,R}$ are real analytic in $R,u,v$. There are
$C<\infty$ and $0<\rho<1$, independent of $j$, such that, writing
$a=|u|+|v|$,
\begin{align}\label{eq:bridge-estimates}
 y_i&=\frac{\sinh((j-i)\chi_R)}{\sinh(j\chi_R)}u
       +\frac{\sinh(i\chi_R)}{\sinh(j\chi_R)}v
       +O\bigl(a^3(\rho^i+\rho^{j-i})\bigr),\\
 |y_i|&\le Ca(\rho^i+\rho^{j-i}),\qquad 0\le i\le j.\notag
\end{align}
The analytic bounds for the bridge and for $W_{j,R}-jg_R$, and
every fixed parameter derivative of these functions, are uniform in $j$. Moreover
\begin{equation}\label{eq:action-jet}
 W_{j,R}(u,v)=jg_R+\tfrac12(u,v)H_{j,R}(u,v)^t+O(a^4),
\end{equation}
with uniform analytic remainder bounds.
\end{lemma}
\begin{proof}
Here $\psi_0=\psi_1=R-\sqrt{R^2-y^2}$ and
$c_0=c_1=c_R$. The endpoint solution of the constant Jacobi equation
is the displayed linear term. Formula~\eqref{eq:v3-green} becomes
\begin{equation}\label{eq:green}
 G_{ik}=g_R\frac{\sinh(i\chi_R)\sinh((j-k)\chi_R)}
                    {\sinh\chi_R\sinh(j\chi_R)},\qquad G_{ki}=G_{ik}
 \quad(1\le i\le k<j).
\end{equation}
For any fixed $e^{-\chi_*}<\rho<1$ and $w_i=\rho^i+\rho^{j-i}$,
\begin{equation}\label{eq:weighted-green}
 \sum_k e^{-\chi_*|i-k|}w_k\le C_\rho w_i.
\end{equation}
Splitting at $k=i$ proves this by the geometric ratios
$e^{-\chi_*}/\rho$ and $e^{-\chi_*}\rho$.
The graphs are even and analytic on a common complex neighborhood
for $R\in K$. The even analytic case of Lemma~\ref{lem:g-relative}
therefore gives the cubic bridge correction, quartic action remainder,
uniqueness throughout the fixed small box and all uniform analytic
bounds asserted here. Its proof applies with interior row margin
$2/R$ and endpoint margin $1/R$, so none of these constants depends
on $j$. For $j=1$ there are no interior variables and
\eqref{eq:ell} gives the same assertions directly.
\end{proof}

\begin{lemma}[Effective Hessian and twist]\label{lem:hessian}
The endpoint Hessian and mixed derivative at the normal segment are
\begin{equation}\label{eq:hessian}
 H_{j,R}=\frac{s_R}{g_R}
 \begin{pmatrix}\coth(j\chi_R)&-\csch(j\chi_R)\\
                -\csch(j\chi_R)&\coth(j\chi_R)\end{pmatrix},
 \qquad
 d^0_{j,R}:=-W_{j,R,uv}(0,0)=\frac{s_R}{g_R\sinh(j\chi_R)}.
\end{equation}
In particular
\begin{equation}\label{eq:hessian-det}
 \sqrt{\det H_{j,R}}=\frac{s_R}{g_R}=\frac1{R\sqrt{g_R}},
 \qquad
 \frac{d^0_{j,R}}{\sqrt{\det H_{j,R}}}=\frac1{\sinh(j\chi_R)}.
\end{equation}
Both eigenvalues of $H_{j,R}$ are bounded above and below by positive
constants independent of $j$ and $R\in K$.
\end{lemma}
\begin{proof}
For the quadratic action the derivative at its first endpoint is
$(c_Ru-y_1^{\rm lin})/g_R$. Insert the bridge formula and use
$c_R\sinh(j\chi_R)-\sinh((j-1)\chi_R)
=s_R\cosh(j\chi_R)$. Reversal gives the other endpoint derivative,
proving \eqref{eq:hessian}. The identity
$\coth^2x-\csch^2x=1$ proves \eqref{eq:hessian-det}.
The eigenvalues are
$(s_R/g_R)\tanh(j\chi_R/2)$ and
$(s_R/g_R)\coth(j\chi_R/2)$. Their asserted bounds follow from
$j\ge1$, $\chi_R\ge\chi_*>0$ and compactness of $K$.
\end{proof}

An absolute bound $W_{uv}=-d^0_{j,R}+O(a^2)$ would be useless for large
$j$. The next estimate is the relevant replacement.

\begin{lemma}[Relative flux determinant]\label{lem:relative}
On a fixed endpoint neighborhood,
\begin{equation}\label{eq:relative}
 -W_{j,R,uv}(u,v)=d^0_{j,R}\,b_{j,R}(u,v),\qquad
 b_{j,R}=1+O(u^2+v^2)>0.
\end{equation}
The functions $b_{j,R}$ are even under $(u,v)\mapsto(-u,-v)$ and
have uniform analytic bounds, including every fixed radius and endpoint
derivative, independently of $j$.
\end{lemma}
\begin{proof}
The even analytic case of Lemma~\ref{lem:g-relative} proves the
relative estimate and every fixed derivative bound. To identify its
exact circular normalization, put
$b_i=-\ell_{R,uv}(y_i,y_{i+1})$ along the stationary segment. Its
corner-cofactor identity is
\begin{equation}\label{eq:twist-product}
 -W_{j,R,uv}=\frac{\prod_{i=0}^{j-1}b_i}{\det H_{\rm int}},
\end{equation}
with empty determinant one when $j=1$. Evenness gives
$\sum_i|\log(g_Rb_i)|\le Ca^2$. If
$E=H_{\rm int}-H^0_{\rm int}$, the same even improvement gives
$\|E\|_1\le Ca^2$ and the exact convergent expansion
\begin{equation}\label{eq:log-det}
 \log\frac{\det H_{\rm int}}{\det H^0_{\rm int}}
 =\sum_{m\ge1}\frac{(-1)^{m+1}}m
   \tr[((H^0_{\rm int})^{-1}E)^m]=O(a^2).
\end{equation}
The trace-norm and differentiated estimates for this expansion were
proved in the general lemma; no factor equal to the matrix dimension
is introduced. Comparing \eqref{eq:twist-product} with its value at
zero and using \eqref{eq:hessian} proves \eqref{eq:relative}.
Uniqueness and simultaneous reversal of the coordinates prove evenness.
Positivity follows from the positive twists and interior Hessian.
\end{proof}

\begin{proposition}[Physical flux in endpoint coordinates]\label{prop:endpoint-flux}
For each of the six oriented nearest-neighbor segments, the small
endpoint coordinates $(u,v)$ parametrize its outgoing initial collision
state one-to-one, and
\begin{equation}\label{eq:endpoint-flux}
 \dd\nu=\frac{-W_{j,R,uv}(u,v)}{4\pi R}\dd u\dd v.
\end{equation}
Moreover $S_j\tau_R=W_{j,R}(u,v)$ on this chart.
\end{proposition}
\begin{proof}
The first variation of length gives
$\partial_uW=-v_0^+\cdot q_0'(u)$. If $s_0$ is boundary arclength,
the outgoing tangential momentum is
$p_0=-W_u/(\dd s_0/\dd u)$. Thus
$|\dd s_0\dd p_0|=|W_{uv}|\dd u\dd v$.
The mixed derivative is nonzero by Lemma~\ref{lem:relative}, so this
is a local coordinate change. At fixed $u$, $p_0$ is strictly monotone
in $v$ throughout the small box, proving injectivity. The stationary
path is physical by Lemma~\ref{lem:localization}; uniqueness identifies
its successive states with the iterates of $B_R$. Normalization of
$\dd s_0\dd p_0$ by its total mass $4\pi R$ proves the measure formula.
The length identity follows by adding the actual flight lengths.
\end{proof}
